\uuid{XeF4}
\titre{Convergence}
\auteur{bodin ; hanani ; mzari}
\organisation{exo7}
\licence{CC BY-NC-SA 4.0}
\datecreate{2020-01-20}
\niveau{L1}
\module{Analyse}
\chapitre{Suite}
\sousChapitre{Autre}
\difficulte{1}
\theme{Suites, convergence}
\exo7qcmid{lille-1.suites.009}

\contenu{
\begin{question}
\qclassification{121.00}
\qtags{motcle=convergence}
Soit $\displaystyle u_n=\ln \left(1+n\mathrm{e}^{-n}\right)$. Quelles sont les bonnes réponses ?
\begin{answers}
    \good{La suite $(u_n)$ est bornée.}
    \bad{$\displaystyle \lim _{n\to +\infty}u_n=+\infty$}
    \good{$\displaystyle \lim _{n\to +\infty}u_n=0$}
    \bad{La suite $(u_n)$ est divergente.}
\end{answers}
\begin{explanations}
Par croissances comparées, $\displaystyle \lim _{n\to +\infty}n\mathrm{e}^{-n}=0$ et, par continuité de la fonction logarithme
$$\lim _{n\to +\infty}u_n=\ln \left[\lim _{n\to +\infty}\left(1+n\mathrm{e}^{-n}\right)\right]=\ln (1)=0.$$
Donc, $(u_n)$ converge et sa limite est $0$. En outre, elle est bornée comme toute suite convergente.
\end{explanations}
\end{question}
}
